\documentclass[11pt,a4paper]{article}

% ─── Packages ────────────────────────────────────────────────────────────────
\usepackage[T1]{fontenc}
\usepackage[utf8]{inputenc}
\usepackage{lmodern}
\usepackage[margin=2.5cm]{geometry}
\usepackage{amsmath,amssymb,amsthm}
\usepackage{bm}
\usepackage{booktabs}
\usepackage{array}
\usepackage{xcolor}
\usepackage{hyperref}
\usepackage{framed}
\usepackage{url}

% ─── Colours ─────────────────────────────────────────────────────────────────
\definecolor{biogreen}{RGB}{30,120,70}
\definecolor{mathblue}{RGB}{10,60,140}
\definecolor{alertred}{RGB}{150,20,20}
\definecolor{warmgrey}{RGB}{80,70,60}
\definecolor{shadecolor}{RGB}{245,248,255}

\hypersetup{colorlinks=true,linkcolor=mathblue,urlcolor=mathblue}

% ─── Theorem-like environments ───────────────────────────────────────────────
\theoremstyle{definition}
\newtheorem{task}{Task}
\newtheorem{question}{Question}

\newenvironment{limitation}[1]{%
  \begin{leftbar}%
  \noindent{\bfseries\color{alertred}Limitation: #1}\\[4pt]%
}{%
  \end{leftbar}%
}

\newenvironment{improvement}[1]{%
  \begin{snugshade}%
  \noindent{\bfseries\color{biogreen}Proposed improvement: #1}\\[4pt]%
}{%
  \end{snugshade}%
}

% ─── Notation (identical to handout) ────────────────────────────────────────
\newcommand{\sig}[2]{S_{#1}(#2)}
\newcommand{\dS}[2]{\Delta S_{#1}(#2)}
\newcommand{\Nb}[2]{\mathcal{N}_{#1}(#2)}
\newcommand{\JE}[2]{\mathcal{J}_{#1}(#2)}
\newcommand{\FJ}[2]{\mathcal{F}_{#1}(#2)}
\newcommand{\expo}[2]{E_{#1}(#2)}
\newcommand{\RD}{\mathrm{RD}}
\newcommand{\RR}{\mathrm{RR}}
\newcommand{\pE}{p_{\mathrm{exp}}}
\newcommand{\pU}{p_{\mathrm{unexp}}}
\newcommand{\1}[1]{\mathbf{1}\!\left[#1\right]}
\newcommand{\Thr}{\theta}

% ─── Utility ─────────────────────────────────────────────────────────────────
\definecolor{pathbg}{RGB}{236,241,252}
\newcommand{\cmd}[1]{\texttt{\small #1}}
\newcommand{\param}[1]{\texttt{#1}}
\newcommand{\file}[1]{\path{#1}}
% \filepath: long paths on their own centred, shaded line (1pt smaller font)
\newcommand{\filepath}[1]{%
  \par\smallskip
  \begin{center}%
    \colorbox{pathbg}{\small\ttfamily\strut\enspace #1\enspace}%
  \end{center}%
  \smallskip\noindent
}

% ─────────────────────────────────────────────────────────────────────────────
\begin{document}

\begin{center}
  {\LARGE\bfseries\color{mathblue}
    Spatiotemporal Signal Propagation\\[4pt]
    Tutorials 2-3}\\[10pt]
  {\large Systems Biology, University of Warsaw}\\[6pt]
  {\small\color{warmgrey}
    Prerequisites: \textit{Spatiotemporal Signal Propagation --- Mathematical
    Formalization} (companion handout).
    Notation is identical throughout.}\\[4pt]
  {\small\color{warmgrey}
    Scripts: \file{scripts/spatiotemporal\_signal\_propagation.py},
    \file{scripts/compare\_spatiotemporal\_behavior.py}}
\end{center}

\vspace{0.5em}
\noindent\rule{\linewidth}{0.8pt}
\vspace{0.5em}

\begin{snugshade}
\noindent\textbf{Session overview.}
\begin{description}
  \item[Part~I (90~min)] Guided exploration of the analysis pipeline.
    You will run single-block and cohort-level analyses, interpret outputs,
    and develop intuition for the effect of parameter choices.
  \item[Part~II (90~min)] Critical analysis and algorithmic improvement.
    You will identify three concrete mathematical limitations of the current
    approach, derive improved estimators on paper, and implement them.
\end{description}
\end{snugshade}

% =============================================================================
\newpage
\section*{Part~I\quad Guided Exploration (90~min)}
\addcontentsline{toc}{section}{Part I --- Guided Exploration}
% =============================================================================

% ─────────────────────────────────────────────────────────────────────────────
\subsection*{Block~1 (25~min)\quad Your first single-block run}
% ─────────────────────────────────────────────────────────────────────────────

\begin{task}[Setup and first run]\label{task:first-run}
Navigate to the project root and run the single-block analysis on
experiment 1, site 1 (a wild-type condition).

\begin{framed}
\begin{verbatim}
python scripts/spatiotemporal_signal_propagation.py \
    --exp-id 1 --site-id 1 \
    --signal-col ERKKTR_ratio \
    --spatial-radius 60 \
    --future-window-frames 3 \
    --jump-quantile 0.9 \
    --output-dir analysis_outputs
\end{verbatim}
\end{framed}

The script prints a short report. Locate each printed quantity in the
mathematical formalism of the companion handout.

\medskip
\noindent\textbf{(a)} How many nodes $|V|$ and spatial edges $|E_{\mathrm{sp}}|$
were created? What is the ratio $|E_{\mathrm{sp}}|/|V|$? What does this ratio
tell you about the local density of the cell monolayer?



\noindent\textbf{(b)} What is the data-driven threshold $\Thr_{0.90}$?
Using Definition~4 from the handout, write out in words what this value means
for a cell at an arbitrary node $(k,t)$.



\noindent\textbf{(c)} Record $\pE$, $\pU$, $\RD$, and $\RR$.
Is the relative risk $\RR$ greater than, equal to, or less than~1?
State in one biological sentence what this implies.


\end{task}

\begin{task}[Inspect the output tables]\label{task:tables}
Open the following file in a spreadsheet or with Python
(\cmd{pd.read\_csv(\ldots, compression='gzip')}):
\filepath{analysis\_outputs/exp\_1\_site\_1\_ERKKTR\_ratio/nodes.csv.gz}

\medskip
\noindent\textbf{(a)} For three tracks of your choice, plot $\sig{k}{t}$ over time
(the column \param{signal\_value}).
Mark the frames where $\JE{k}{t}=1$ (\param{jump\_event}~$=$ True) with a
vertical line or a dot.
Are the jumps isolated single frames, or do they appear in bursts?



\noindent\textbf{(b)} Filter to rows where \param{future\_self\_jump} is True
\emph{and} \param{neighbor\_jump\_now} is True.
How many such rows are there?
Express this count as a fraction of all \param{neighbor\_jump\_now} = True rows.
Verify that this matches $\pE$ from Task~\ref{task:first-run}.



\noindent\textbf{(c)} Compute the mean and standard deviation of
\param{neighbor\_count} across all nodes.
Sketch (by hand or code) a histogram of \param{neighbor\_count}.
At radius $r = 60$, is the neighbourhood sparse or dense?


\end{task}

% ─────────────────────────────────────────────────────────────────────────────
\subsection*{Block~2 (30~min)\quad Parameter sensitivity}
% ─────────────────────────────────────────────────────────────────────────────

Keeping the same block (exp~1, site~1, ERK signal), re-run the script with the
following parameter variations.
Fill in the table below each time; record at minimum $\Thr$, $|E_{\mathrm{sp}}|$,
$\pE$, $\pU$, $\RR$.

\begin{task}[Varying the spatial radius $r$]\label{task:radius}
Run with $r \in \{30,\; 60,\; 90,\; 150\}$.

\begin{center}
\renewcommand{\arraystretch}{1.4}
\begin{tabular}{ccccc}
\toprule
$r$ & $|E_{\mathrm{sp}}|$ & $\pE$ & $\pU$ & $\RR$ \\
\midrule
30  & & & & \\
60  & & & & \\
90  & & & & \\
150 & & & & \\
\bottomrule
\end{tabular}
\end{center}

\medskip
\noindent\textbf{(a)} As $r$ grows, $|E_{\mathrm{sp}}|$ grows roughly how fast?
(Hint: for a uniform cell density $\rho$ [cells per unit area], the expected
number of neighbours within radius $r$ is $\mathbb{E}[n(k,t)] \approx \pi r^2 \rho$.
Does $|E_{\mathrm{sp}}|$ scale as $r^2$?)



\noindent\textbf{(b)} Does $\RR$ increase or decrease with $r$?
Explain why a very large $r$ should push $\RR$ toward~1.
(Hint: consider what happens when almost every cell is a neighbour of every
other cell.)


\end{task}

\begin{task}[Varying the future window $W$]\label{task:window}
Run with $W \in \{1,\; 3,\; 6,\; 12\}$ (keeping $r=60$).

\begin{center}
\renewcommand{\arraystretch}{1.4}
\begin{tabular}{ccccc}
\toprule
$W$ (frames) & $W \cdot 5$ min & $\pE$ & $\pU$ & $\RR$ \\
\midrule
1  & 5 min  & & & \\
3  & 15 min & & & \\
6  & 30 min & & & \\
12 & 60 min & & & \\
\bottomrule
\end{tabular}
\end{center}

\medskip
\noindent\textbf{(a)} Does $\pE$ increase or decrease with $W$?
Is the same trend observed for $\pU$?
Why does both rates rising not necessarily imply that $\RR$ changes?



\noindent\textbf{(b)} At what window length does $\RR$ peak?
What does this suggest about the timescale of cell-to-cell ERK coordination
in this dataset?


\end{task}

\begin{task}[Varying the jump quantile $q$]\label{task:quantile}
Run with $q \in \{0.70,\; 0.80,\; 0.90,\; 0.95\}$ (keeping $r=60$, $W=3$).

\begin{center}
\renewcommand{\arraystretch}{1.4}
\begin{tabular}{ccccc}
\toprule
$q$ & $\Thr_q$ & $\pE$ & $\pU$ & $\RR$ \\
\midrule
0.70 & & & & \\
0.80 & & & & \\
0.90 & & & & \\
0.95 & & & & \\
\bottomrule
\end{tabular}
\end{center}

\medskip
\noindent\textbf{(a)} As $q$ increases, fewer nodes satisfy $\JE{k}{t}=1$.
How does the fraction of jump events scale with $q$?



\noindent\textbf{(b)} Does increasing $q$ (i.e.\ using a stricter threshold
for calling a ``jump'') increase or decrease $\RR$?
Propose a biological interpretation: are rarer, larger jumps more or less
spatially coordinated in this dataset?


\end{task}

% ─────────────────────────────────────────────────────────────────────────────
\subsection*{Block~3 (35~min)\quad Comparing individual conditions and
  running a cohort}
% ─────────────────────────────────────────────────────────────────────────────

\begin{task}[Manual pairwise comparison]\label{task:pairwise}
According to the metadata file
\filepath{01-readme-experiment-description\_2022-04-05.csv}
sites 1--4 are WT and sites 9--12 are PIK3CA E545K.
Run the single-block script separately for \textbf{site~1 (WT)} and
\textbf{site~9 (PIK3CA E545K)}.

\medskip
\noindent Fill in the table (use the same $r=60$, $W=3$, $q=0.90$):

\begin{center}
\renewcommand{\arraystretch}{1.4}
\begin{tabular}{lccccc}
\toprule
Condition & Site & $\Thr$ & $\pE$ & $\pU$ & $\RR$ \\
\midrule
WT           & 1 & & & & \\
PIK3CA E545K & 9 & & & & \\
\bottomrule
\end{tabular}
\end{center}

\medskip
\noindent\textbf{(a)} Is $\RR$ higher or lower for the PIK3CA mutant compared
to WT?
PIK3CA E545K is a gain-of-function oncogenic mutation that hyperactivates PI3K
signalling. Based on your result, does hyperactive PI3K increase or decrease
ERK-activity coordination among neighbours?



\noindent\textbf{(b)} Can you directly compare the two $\Thr$ values?
Why or why not?
(This question anticipates Limitation~L2 in Part~II.)


\end{task}

\begin{task}[Cohort-level analysis with the comparison script]\label{task:cohort}
Run the multi-block comparison script to aggregate all available
blocks and compare by mutation:

\begin{framed}
\begin{verbatim}
python scripts/compare_spatiotemporal_behavior.py \
    --signal-col ERKKTR_ratio \
    --group-by mutation \
    --spatial-radius 60 \
    --future-window-frames 3 \
    --jump-quantile 0.9 \
    --exclude-mutations ErbB2 \
    --output-dir analysis_outputs
\end{verbatim}
\end{framed}

Open the group-level summary:
\filepath{analysis\_outputs/comparison\_mutation\_ERKKTR\_ratio/group\_level\_summary.csv}

\medskip
\noindent\textbf{(a)} Which mutation has the \emph{highest} $\overline{\RR}$?
Which has the lowest?
Does the ordering match your expectation from the PI3K pathway hierarchy
(PI3K $\to$ AKT $\to$ PTEN)?



\noindent\textbf{(b)} Note the column \param{n\_blocks} for each group.
Some mutations have more replicate blocks than others.
Why is $\overline{\RR}$ (the mean over blocks) a fragile summary when
\param{n\_blocks} is small?



\noindent\textbf{(c)} Repeat the cohort run with \param{--signal-col FoxO3A\_ratio}.
Fill in the comparison table:

\begin{center}
\renewcommand{\arraystretch}{1.4}
\begin{tabular}{lcc}
\toprule
Mutation & $\overline{\RR}$ (ERK) & $\overline{\RR}$ (FoxO3A) \\
\midrule
WT           & & \\
AKT1 E17K    & & \\
PIK3CA E545K & & \\
PIK3CA H1047R& & \\
PTEN del     & & \\
\bottomrule
\end{tabular}
\end{center}

Is ERK or FoxO3A more spatially coordinated overall?
Given that FoxO3A is a direct AKT substrate while ERK is a MAPK downstream of
RAS, does the difference make biological sense?


\end{task}

% =============================================================================
\newpage
\section*{Part~II\quad Critical Analysis and Algorithmic Improvement (90~min)}
\addcontentsline{toc}{section}{Part II --- Critical Analysis and Algorithmic Improvement}
% =============================================================================

\noindent The analysis you have run so far is a solid first pass.
In this part you will identify three mathematical limitations, see concretely
why each one matters, and then derive and implement an improved estimator.
All improvements are achievable with NumPy/pandas without modifying the
original scripts beyond a new analysis function.

% ─────────────────────────────────────────────────────────────────────────────
\subsection*{Limitation~L1 (25~min)\quad Simultaneity: exposure and response
  are not properly ordered in time}
% ─────────────────────────────────────────────────────────────────────────────

\subsubsection*{The problem}

Recall the exposure indicator (Definition~5 of the handout):
\[
  \expo{k}{t} \;=\; \1{m(k,t) > 0}
  \qquad\text{where}\qquad
  m(k,t) = \sum_{(k',t)\,\in\,\Nb{k}{t}} \JE{k'}{t}.
\]
The neighbour's jump $\JE{k'}{t}$ is measured \emph{at the same frame $t$}
as the focal cell is observed.
But the future self-jump is defined as
$\FJ{k}{t} = \1{\exists\, t' \in (t, t{+}W] : \JE{k}{t'}=1}$.

\medskip
\noindent\textbf{The temporal ordering problem.}
Both the neighbour's jump $\JE{k'}{t}$ and the future-jump window
$(t, t{+}W]$ are anchored to the same frame $t$.
If cell $k'$ jumps at frame $t$ and cell $k$ jumps at frame $t+1$, the
current analysis classifies this as a ``successful propagation event''.
But it could equally be that both cells respond to the same upstream stimulus
arriving at frame~$t$ --- there is no \emph{temporal precedence} in the
measurement.
More critically, the current formula also counts cases where \emph{cell $k$
itself jumps simultaneously} with its neighbour (at frame $t$) and then again
within $(t, t+W]$, even if the two events are independent.

\begin{center}
\renewcommand{\arraystretch}{1.4}
\begin{tabular}{ccccc}
\toprule
Frame & $\JE{k}{t}$ & $\JE{k'}{t}$ & $\expo{k}{t}$ & $\FJ{k}{t}$
  (W=3) \\
\midrule
$t$   & 0 & 1 & 1 & ? \\
$t+1$ & 1 & 0 & 0 & \\
$t+2$ & 0 & 0 & 0 & \\
$t+3$ & 0 & 0 & 0 & \\
\bottomrule
\end{tabular}
\end{center}

\begin{question}
In the table above, what is $\FJ{k}{t}$?
Does the sequence truly suggest that cell $k'$ ``caused'' cell $k$ to jump?
Can you construct a scenario where the assignment $\FJ{k}{t}=1$ is misleading?

\end{question}

\subsubsection*{Proposed improvement: lagged exposure}

\begin{improvement}{Lagged exposure $E_k^{(\tau)}(t)$}
For a \emph{lag} $\tau \in \{1, 2, \ldots\}$ (in frames), define the
\textbf{lagged exposure indicator}
\[
  \expo{k}{t}^{(\tau)}
  \;=\; \1{m(k,\,t-\tau) > 0},
\]
where $m(k, t-\tau)$ counts how many of cell $k$'s spatial neighbours at
frame $t-\tau$ were jumping at that earlier frame.
The lagged relative risk is then
\[
  \RR^{(\tau)}
  \;=\;
  \frac{\Pr\!\bigl(\FJ{k}{t}=1 \;\big|\; \expo{k}{t}^{(\tau)}=1\bigr)}
       {\Pr\!\bigl(\FJ{k}{t}=1 \;\big|\; \expo{k}{t}^{(\tau)}=0\bigr)}.
\]
For $\tau = 0$ this reduces to the original $\RR$.
For $\tau \geq 1$, the neighbour's jump \emph{strictly precedes} the focal
cell's look-ahead window, providing genuine temporal precedence.
\end{improvement}

\medskip
\noindent The key insight is: if signal truly propagates from a jumping neighbour
\emph{to} the focal cell, we expect $\RR^{(\tau)} > 1$ for small lags
$\tau \approx 1$--$3$ and then $\RR^{(\tau)} \to 1$ as $\tau$ grows beyond the
signalling timescale.

\begin{task}[Implement and plot $\RR^{(\tau)}$ vs.\ $\tau$]\label{task:lag}
Load the node table you saved in Task~\ref{task:tables}
(\file{nodes.csv.gz} for exp~1, site~1).

\medskip
\noindent\textbf{Step 1.}
The current table already contains \param{neighbor\_jump\_now} (which is
$\expo{k}{t}^{(0)}$).
You need to construct $\expo{k}{t}^{(\tau)}$ for $\tau = 1, 2, 3, 4, 5$.
Within each track, shift the column \param{neighbor\_jump\_now} forward
by $\tau$ frames:
\begin{framed}
\begin{verbatim}
# Example in pandas (nodes_df loaded from nodes.csv.gz)
nodes_df = nodes_df.sort_values(['track_id', 'Image_Metadata_T'])
for tau in range(1, 6):
    nodes_df[f'lagged_exposure_{tau}'] = (
        nodes_df.groupby('track_id')['neighbor_jump_now']
        .shift(tau)          # shift forward by tau frames within track
        .fillna(False)
    )
\end{verbatim}
\end{framed}

\noindent\textbf{Step 2.}
For each $\tau \in \{0, 1, 2, 3, 4, 5\}$, compute $\RR^{(\tau)}$:

\begin{framed}
\begin{verbatim}
import numpy as np
results = {}
for tau in range(6):
    col = 'neighbor_jump_now' if tau == 0 else f'lagged_exposure_{tau}'
    exposed   = nodes_df[nodes_df[col] == True]['future_self_jump']
    unexposed = nodes_df[nodes_df[col] == False]['future_self_jump']
    p_exp   = exposed.mean()
    p_unexp = unexposed.mean()
    results[tau] = {'p_exp': p_exp, 'p_unexp': p_unexp,
                    'RR': p_exp / p_unexp if p_unexp > 0 else np.nan}
\end{verbatim}
\end{framed}

\noindent\textbf{Step 3.}
Plot $\RR^{(\tau)}$ against $\tau$ (x-axis: lag in frames; optionally
convert to minutes using $5$ min/frame).
Add a horizontal dashed line at $\RR = 1$ as a reference.

\medskip
\noindent\textbf{(a)} At which lag $\tau^*$ is $\RR^{(\tau)}$ highest?
What is the corresponding physical time $\tau^* \times 5$~min?
Does this suggest a rapid or slow cell-to-cell ERK relay?



\noindent\textbf{(b)} Does $\RR^{(\tau)}$ drop below~1 at large $\tau$,
or does it plateau?
What does each scenario imply about the persistence of spatial ERK
coordination?



\noindent\textbf{(c)} Compare $\RR^{(0)}$ (original) with $\RR^{(1)}$.
Is the original estimator inflated, deflated, or unaffected by the
simultaneity issue for this dataset?


\end{task}

% ─────────────────────────────────────────────────────────────────────────────
\subsection*{Limitation~L2 (30~min)\quad Non-comparable thresholds across blocks}
% ─────────────────────────────────────────────────────────────────────────────

\subsubsection*{The problem}

Recall Definition~4 of the handout: when no global threshold is specified,
each block $b$ uses its own data-driven threshold
$\Thr_b = Q_q(\mathcal{D}_b^+)$,
where $\mathcal{D}_b^+$ is the set of positive increments \emph{within block $b$}.

This creates a serious comparability problem.
Let $\mu_b = \mathbb{E}[\dS{k}{t}^+]$ be the mean positive increment in
block $b$.
If a mutant cell line has systematically larger ERK pulses
($\mu_{\mathrm{mut}} > \mu_{\mathrm{WT}}$), its per-block threshold
$\Thr_{\mathrm{mut}}$ will also be larger.
The consequence:

\medskip
\noindent\textbf{Worked numerical example.}
Suppose we observe (in arbitrary ratio units):

\begin{center}
\renewcommand{\arraystretch}{1.3}
\begin{tabular}{lccc}
\toprule
Block & Positive increments $\mathcal{D}^+$ (sample) & $\Thr_{0.90}$ & \% nodes called jumps \\
\midrule
WT           & $\{0.01,\ldots,0.10\}$, mean $0.03$ & $0.07$ & $10\%$ \\
PTEN del     & $\{0.03,\ldots,0.30\}$, mean $0.09$ & $0.21$ & $10\%$ \\
\bottomrule
\end{tabular}
\end{center}

Both blocks end up labelling exactly $10\%$ of nodes as jumps (by
construction of the quantile threshold).
But PTEN-deleted cells show a \emph{three-fold larger} typical increment.
The per-block threshold hides this biological difference: we are comparing
``top-10\% ERK steps within WT'' to ``top-10\% ERK steps within PTEN del'',
which are intrinsically different events.

\begin{question}
In the above example, if we instead used a single global threshold
$\Thr_{\mathrm{global}} = Q_{0.90}(\mathcal{D}_{\mathrm{all}}^+)$
computed from \emph{all blocks together}, which mutation would have more
jump events?
Does the relative $\RR$ estimate between mutations change?
Explain why this matters for a \emph{cross-group comparison}.

\end{question}

\subsubsection*{Proposed improvement: global cross-block threshold}

\begin{improvement}{Global threshold $\Thr^*$}
Collect all positive signal increments from every block in the comparison:
\[
  \mathcal{D}_{\mathrm{all}}^+
  \;=\; \bigcup_{b\,\in\,\mathcal{B}} \mathcal{D}_b^+
  \;=\; \bigcup_{b\,\in\,\mathcal{B}}
    \bigl\{\dS{k}{t} \;\big|\; (k,t) \in V_b,\; \dS{k}{t} > 0\bigr\}.
\]
Define the \textbf{global threshold}
\[
  \Thr^* = Q_q\!\bigl(\mathcal{D}_{\mathrm{all}}^+\bigr).
\]
Re-run the single-block analysis for each block with this fixed $\Thr^*$
(using the \param{--jump-threshold} flag instead of \param{--jump-quantile}).
Under the global threshold, the same increment size counts as a jump in
every block, making cross-group $\RR$ values directly comparable.
\end{improvement}

\begin{task}[Implement global-threshold comparison]\label{task:global-threshold}
You will need to do this in two steps.

\medskip
\noindent\textbf{Step 1: Estimate $\Thr^*$.}
For each mutation's site(s), load the node table(s) and pool the positive
increments. (If you only have WT data from Task~\ref{task:first-run},
use those.)
\begin{framed}
\begin{verbatim}
import pandas as pd, numpy as np

all_deltas = []
for node_file in your_node_files:          # list of nodes.csv.gz paths
    df = pd.read_csv(node_file)
    pos = df['signal_delta'].dropna()
    all_deltas.append(pos[pos > 0].values)

all_deltas = np.concatenate(all_deltas)
theta_star = float(np.quantile(all_deltas, 0.90))
print(f"Global threshold theta* = {theta_star:.4f}")
\end{verbatim}
\end{framed}

\noindent\textbf{Step 2: Re-run with global threshold.}
\begin{framed}
\begin{verbatim}
python scripts/spatiotemporal_signal_propagation.py \
    --exp-id 1 --site-id 9 \
    --signal-col ERKKTR_ratio \
    --jump-threshold <theta_star> \
    --output-dir analysis_outputs_global
\end{verbatim}
\end{framed}

\noindent\textbf{(a)} Compare the fraction of nodes labelled as jumps between WT
(site~1) and PIK3CA E545K (site~9) when using (i)~per-block thresholds
vs.\ (ii)~the global threshold $\Thr^*$.
Do the two conditions differ in jump frequency under the global threshold?



\noindent\textbf{(b)} Compare $\RR_{\mathrm{WT}}$ and $\RR_{\mathrm{PIK3CA}}$
under per-block vs.\ global thresholds.
Does the ranking of conditions change?
Which comparison do you trust more for a biological claim about
``whose cells are more spatially coordinated''?



\noindent\textbf{(c)} Formally, define the \emph{bias} of the per-block estimator.
Let $\RR_b^{\mathrm{per}}$ be computed with $\Thr_b$ and
$\RR_b^{\mathrm{global}}$ be computed with $\Thr^*$.
What conditions on $\Thr_b$ vs.\ $\Thr^*$ would make
$\RR_b^{\mathrm{per}} > \RR_b^{\mathrm{global}}$?
(Hint: think about which nodes shift from ``not a jump'' to ``a jump''
when the threshold is lowered.)


\end{task}

% ─────────────────────────────────────────────────────────────────────────────
\subsection*{Limitation~L3 (35~min)\quad Binary exposure discards
  neighbourhood-intensity information}
% ─────────────────────────────────────────────────────────────────────────────

\subsubsection*{The problem}

The exposure indicator
\[
  \expo{k}{t} = \1{m(k,t) > 0}
\]
is binary: a node is ``exposed'' even if only \emph{one} of its $n(k,t)$
neighbours is jumping, and it receives the same label whether $1$ or $8$
neighbours are jumping simultaneously.

This collapses the entire neighbourhood into a single bit and discards two
sources of potentially important information:
\begin{enumerate}
  \item \textbf{Fraction of jumping neighbours.}
        A cell surrounded by $8$ jumping neighbours might be more strongly
        stimulated than one with only $1$ jumping neighbour.
  \item \textbf{Intensity of the neighbourhood signal.}
        The mean signal $\bar{S}(k,t)$ of the neighbourhood (already computed
        in the node table as \param{neighbor\_mean\_signal}) carries information
        about the ambient signalling environment that is ignored by the binary
        indicator.
\end{enumerate}

\begin{question}
Consider two nodes $(k,t)$ and $(k',t')$, both with
$n = 8$ spatial neighbours.
Cell $k$ has $m=1$ jumping neighbour; cell $k'$ has $m=8$ jumping neighbours.
Under the current formulation, both are ``exposed'' equally
($\expo{k}{t} = \expo{k'}{t'} = 1$).

Does it seem reasonable to you that both cells should have the
same expected future-jump probability $\pE$?
What biological mechanism would justify equal treatment?
What mechanism would predict that $k'$ (with 8 jumping neighbours) is more
likely to jump in the near future?


\end{question}

\subsubsection*{Proposed improvement A: fractional exposure}

\begin{improvement}{Fractional exposure $\tilde{E}_k(t)$}
Replace the binary indicator with the \textbf{fraction of jumping
neighbours}:
\[
  \tilde{E}_k(t)
  \;=\;
  \begin{cases}
    m(k,t) / n(k,t) & \text{if } n(k,t) > 0, \\
    0               & \text{if } n(k,t) = 0.
  \end{cases}
\]
$\tilde{E}_k(t) \in [0,1]$.
Partition nodes into exposure \emph{bins}
$[0, \frac{1}{B}), [\frac{1}{B}, \frac{2}{B}), \ldots, [\frac{B-1}{B}, 1]$
and compute a separate future-jump rate $p_b$ for each bin.
A dose--response relationship ($p_b$ increasing with bin index) would support
the hypothesis that more jumping neighbours produce stronger induction.
\end{improvement}

\subsubsection*{Proposed improvement B: signal-level exposure}

\begin{improvement}{Neighbourhood-signal exposure $\hat{E}_k(t)$}
Use the \textbf{mean neighbour signal} as a continuous exposure variable:
\[
  \hat{E}_k(t) \;=\; \bar{S}(k,t)
  \;=\; \frac{1}{n(k,t)}\sum_{(k',t)\,\in\,\Nb{k}{t}} \sig{k'}{t}.
\]
Discretize $\hat{E}_k(t)$ into quartiles $Q_1$--$Q_4$ across all nodes with
$n(k,t) > 0$.
Compute the future-jump rate $p_{Q_j}$ for each quartile.
An increasing pattern (higher ambient neighbourhood signal $\Rightarrow$
higher future-jump probability) would indicate that the mean state of the
neighbourhood, not just its ``jumpiness'', predicts self-activation.
\end{improvement}

\begin{task}[Implement dose--response analysis]\label{task:dose-response}
Using the node table from Task~\ref{task:tables}, implement both improvements.

\medskip
\noindent\textbf{Step 1: Fractional exposure.}
\begin{framed}
\begin{verbatim}
df = pd.read_csv('nodes.csv.gz')

# fractional exposure: use neighbour_jump_count / neighbour_count
df['frac_exposure'] = (
    df['neighbor_jump_count'] / df['neighbor_count'].replace(0, np.nan)
).fillna(0.0)

# bin into 5 levels: 0, (0,0.25], (0.25,0.5], (0.5,0.75], (0.75,1]
bins   = [-0.001, 0.0, 0.25, 0.50, 0.75, 1.01]
labels = ['0', '(0,0.25]', '(0.25,0.5]', '(0.5,0.75]', '(0.75,1]']
df['frac_bin'] = pd.cut(df['frac_exposure'], bins=bins, labels=labels)

dose_response = (df.groupby('frac_bin', observed=True)['future_self_jump']
                   .agg(['mean', 'count'])
                   .rename(columns={'mean': 'future_jump_rate',
                                    'count': 'n_nodes'}))
print(dose_response)
\end{verbatim}
\end{framed}

\noindent\textbf{Step 2: Neighbourhood-signal exposure.}
\begin{framed}
\begin{verbatim}
has_nb = df['neighbor_count'] > 0
df.loc[has_nb, 'nb_signal_quartile'] = pd.qcut(
    df.loc[has_nb, 'neighbor_mean_signal'], q=4,
    labels=['Q1','Q2','Q3','Q4']
)
nb_dose = (df.groupby('nb_signal_quartile', observed=True)['future_self_jump']
             .agg(['mean','count'])
             .rename(columns={'mean':'future_jump_rate','count':'n_nodes'}))
print(nb_dose)
\end{verbatim}
\end{framed}

\noindent\textbf{(a)} Plot the dose--response curves for both improvements
(bar chart or line plot: x-axis = bin/quartile, y-axis = $p_b$).
Does the fractional-exposure curve ($\tilde{E}$) show a monotone
dose--response?
What would a flat curve (all bins equal) imply biologically?



\noindent\textbf{(b)} Now compute a single summary statistic analogous to $\RR$
for the fractional exposure.
Define the \emph{graded relative risk}:
\[
  \widetilde{\RR}
  \;=\;
  \frac{p_{B}}{p_{0}},
\]
where $p_B$ is the future-jump rate in the highest fraction bin and $p_0$ is
the rate in the zero-exposure bin ($\tilde{E}=0$).
Compare $\widetilde{\RR}$ with the original $\RR$ from Task~\ref{task:first-run}.
Which estimator is larger?
Why does $\widetilde{\RR}$ represent a ``cleaner'' signal?



\noindent\textbf{(c)} Does the neighbourhood signal level $\hat{E}_k(t)$
independently predict future self-jumping (beyond the binary jump indicator)?
In other words, do cells in high-signal neighbourhoods (Q4) jump more in the
future even when $\expo{k}{t}=0$ (no jumping neighbour)?
Filter the data to \param{neighbor\_jump\_now}~$=$~False and repeat
the quartile analysis for this sub-population.


\end{task}

% ─────────────────────────────────────────────────────────────────────────────
\subsection*{Synthesis (bonus, 10~min)}
% ─────────────────────────────────────────────────────────────────────────────

\begin{task}[Combining all three improvements]\label{task:synthesis}
You now have three independently derived improvements:

\begin{center}
\renewcommand{\arraystretch}{1.3}
\begin{tabular}{lll}
\toprule
Limitation & Improvement & Key new quantity \\
\midrule
L1: simultaneity & Lagged exposure & $\expo{k}{t}^{(\tau)}$, $\RR^{(\tau)}$ \\
L2: non-comparable thresholds & Global threshold & $\Thr^*$, $\RR_b^{\mathrm{global}}$ \\
L3: binary exposure & Fractional/signal exposure & $\tilde{E}_k(t)$, $\widetilde{\RR}$ \\
\bottomrule
\end{tabular}
\end{center}

Suppose you wanted to combine improvements L1 and L3 into a single estimator.
Write down the mathematical definition of the \emph{lagged fractional exposure}
\[
  \tilde{E}_k^{(\tau)}(t) \;=\; \frac{m(k,\,t-\tau)}{n(k,\,t-\tau)},
\]
and the corresponding graded relative risk
\[
  \widetilde{\RR}^{(\tau)} \;=\; \frac{p_B^{(\tau)}}{p_0^{(\tau)}},
\]
where $p_B^{(\tau)}$ is the future-jump rate for nodes in the highest
fractional-exposure bin \emph{at lag $\tau$}.

\medskip
\noindent\textbf{(a)} Implement $\widetilde{\RR}^{(\tau)}$ for
$\tau \in \{0, 1, 2, 3, 4\}$ and plot the resulting lag curve.
Does it peak at the same $\tau^*$ as the binary $\RR^{(\tau)}$ from
Task~\ref{task:lag}?



\noindent\textbf{(b)} \textbf{Open discussion (no single correct answer):}
What further mathematical improvement would you propose?
Some options to consider:
\begin{itemize}
  \item Accounting for the directionality of cell movement (use a
        directional neighbourhood cone rather than a symmetric ball).
  \item Using a continuous signal increment $\dS{k}{t}$ instead of the
        binary $\JE{k}{t}$ to define a ``soft jump'' probability via a
        sigmoid function.
  \item Replacing the empirical mean estimators for $\pE$ and $\pU$ with a
        mixed-effects logistic regression that treats the track as a random
        effect, properly handling the non-independence of nodes.
\end{itemize}
Describe your chosen improvement in mathematical notation and sketch
the implementation.



\end{task}

% ─────────────────────────────────────────────────────────────────────────────
\bigskip
\noindent\rule{\linewidth}{0.8pt}
\begin{center}
{\small\color{warmgrey}
  Notation follows the companion handout exactly.
  All tasks assume Python~3, pandas, NumPy, and the project virtual environment.
  Refer to the notation table in the Appendix of the handout for symbol
  definitions.}
\end{center}

\end{document}
